\documentclass[uplatex]{jsreport}
\usepackage{docmute}
\usepackage{myStyle}
\label{chapter24}
\setcounter{chapter}{23-1}

\begin{document}

\chapter{確率分布とデータの関係;尤度}
検定では平均因果効果を検証するため，平均値や(不偏)分散などの標本統計量がどのような確率分布に従うかを参照し，結果を判定するという使い方を行なってきました。検定はこのように，母数を推定する方法に，結果を判定するロジックが組み合わさったものなのでした。この時の母数の推定方法は，標本平均や不偏分散を推定値とする方法で，モーメント法とよばれるものです。これに対し，モデルである確率分布をデータに最も近しい形に調節することで，そのパラメータを推定値とする方法があります。ここではこの方法，最尤推定法とその根幹にある尤度という考え方を解説します。
\section{確率関数と尤度}
\section{最尤推定}
\section{対数と尤度}
\section{正規分布の尤度関数}
\section{課題}

\end{document}
